% Final project report, Web Data Integration HWS 2026
% Springer LNCS template: llncs.cls v2.26. Build: latexmk -pdf report.tex
%
% Rules from the course slides (slide 15):
%   - 12 pages, counted without title page, table of contents and literature list;
%     every extra page (appendix included) lowers the mark by 0.33
%   - due Sunday, 29 November 2026, 23:59, via ILIAS, together with the code and (a subset of) the data
%   - cite sources properly, preferred style [Author, year]
%   - hand in a table on who did what together with the report (slide 7): contributions.md
% The result tables in tables/ are written by scripts/make_report_tables.py.
\documentclass[runningheads,a4paper]{llncs}
\usepackage[T1]{fontenc}
\usepackage{graphicx}
\usepackage{booktabs}
\usepackage{tabularx}
\usepackage[square,authoryear]{natbib}
\usepackage[hidelinks]{hyperref}
\usepackage{xurl}
\urlstyle{same}
\renewcommand{\bibsection}{\section*{References}}
\renewcommand{\bibfont}{\small}
\newcommand{\mv}{\,(MV)}
\newcolumntype{L}{>{\raggedright\arraybackslash}X}
\newcolumntype{P}[1]{>{\raggedright\arraybackslash}p{#1}}
\setlength{\textfloatsep}{8pt plus 2pt minus 2pt}
\setlength{\floatsep}{6pt plus 2pt minus 2pt}
\setlength{\intextsep}{6pt plus 2pt minus 2pt}

\begin{document}

\title{Integrating Worldwide University Data from Wikidata, ROR and DEQAR}
\subtitle{Project Report -- Web Data Integration, HWS 2026}
\titlerunning{Integrating Worldwide University Data}

\author{Laith Sandouk \and Mohammad Hadad \and Sanjeevani Rajpurohit \and Chellapuram Keerthan Reddy
\and Lakpa Dolma Sherpa}
\authorrunning{Web Data Integration, HWS 2026}
\institute{University of Mannheim, Germany}

\maketitle

\begin{abstract}
We integrate three open sources on universities: Wikidata, the Research Organization Registry (ROR)
and the Database of External Quality Assurance Results (DEQAR). The 61,807 records describe 44,340
distinct institutions, 14,630 of them in at least two sources and 2,837 in all three. We translate the
sources into one XML schema with 21 attributes, build a gold standard from the identifier links of the
sources and from hand-picked corner cases, compare blocking strategies, similarity measures, rules
and learned matchers with PyDI, and fuse the matched records with seven conflict resolution
strategies. Learned matchers reach an F1 of 96--99\% on the gold standard, 4--8 points above the
best matcher that uses names only; the fused dataset fills attributes such as the founding year for 99\%
of the entities that are in at least two sources.
\end{abstract}

\section{Introduction}

Students, researchers and policy makers who compare universities worldwide meet registries,
knowledge bases and quality assurance databases that differ in format, definitions and coverage.
ROR has reliable names and locations but few descriptive facts, Wikidata adds founding years,
student numbers and memberships but is incomplete, and DEQAR records mergers, closures and external
quality assurance reports, mostly of European institutions. Our goal is one dataset that describes
every university exactly once, with the best value for each attribute.

As in the project abstract, a \emph{university} is a higher education institution that awards
academic degrees and is an organisation of its own; a college or institute that meets both
conditions counts, a branch, faculty or campus does not. The sources draw this line differently
(ROR also holds branches, DEQAR faculties and closed predecessors), and none of them records
whether an institution awards degrees. We therefore keep such records in the input and use them as
hard non-matches of the gold standard (Sect.~\ref{sec:gold}) instead of filtering them by rules the
data cannot support. The integration follows the three phases of the course: data collection and
translation (Sect.~\ref{sec:phase1}), identity resolution (Sect.~\ref{sec:phase2}) and data fusion
(Sect.~\ref{sec:phase3}). All code, the translated data, the gold standards and the fused dataset
are part of the submission (folders \url{scripts/} and \url{data/}).

\section{Phase I: Data Collection and Translation}\label{sec:phase1}

\paragraph{Datasets.}
Our Python scripts retrieved Wikidata and ROR on 2 October 2026 and DEQAR on 9 October 2026 and store
each source as one CSV file with one row per institution; attributes with several values are lists.
\emph{Wikidata} \citep{vrandecic2014wikidata} comes from SPARQL queries against
\url{query.wikidata.org/sparql}: all items that are an instance of \emph{university} (Q3918) or one of
its subclasses, with identifiers replaced by English labels. \emph{ROR} \citep{ror2026dump} is the data
dump v2.13 of 22 September 2026; of its 141,528 records we kept the 30,131 of type \emph{education}
that are not withdrawn. \emph{DEQAR} \citep{eqar2026deqar} is the list of higher education
institutions that EQAR publishes as a nightly CSV file; we merged rows that repeat an institution
once per location and took the country from the ETER identifier where the file has none. The three
sources are edited by volunteers (Wikidata), a curation team (ROR) and quality assurance agencies
(DEQAR). On linked records the founding year of DEQAR equals that of ROR in only 73\% and that of
Wikidata in 74\% of the cases, so the sources are sufficiently distinct to vote against each other.
Table~\ref{tab:datasets} profiles them.

\begin{table}[t]
\caption{Datasets}\label{tab:datasets}
\scriptsize
\setlength{\tabcolsep}{3pt}
\begin{tabularx}{\textwidth}{@{}lP{2.3cm}lP{0.9cm}P{1.1cm}L@{}}
\toprule
Dataset & Source (*) & Format & \# entities & \# attributes & List of attributes (**)\\
\midrule
Wikidata & \url{query.wikidata.org/sparql} & CSV & 22,697 & 16 &
Wikidata id, name, alt.\ names\mv, acronyms\mv, country, country code, city, website, founding
year, latitude\mv, longitude\mv, students\mv, types, member of\mv, ROR id\mv, Wikipedia title\mv\\
\addlinespace[2pt]
ROR & \url{doi.org/10.5281/zenodo.22902037} & CSV & 30,131 & 18 &
ROR id, name, alt.\ names, acronyms\mv, country, country code, region, city, website, domains,
founding year, latitude, longitude, types, status, parent organisation\mv, Wikidata id,
Wikipedia title\mv\\
\addlinespace[2pt]
DEQAR & \url{eqar.eu/qa-results/download-data-sets} & CSV & 8,979 & 13 &
DEQAR id, name, alt.\ names\mv, country, city, website, founding year\mv, closure year\mv,
parent institutions\mv, history\mv, QA reports\mv, last QA report\mv, ROR id\mv\\
\bottomrule
\end{tabularx}

\smallskip
\scriptsize (*) Retrieved with our own scripts (\url{scripts/fetch_*.py}, \url{build_ror_education.py}). (**) (MV) =
more than 30\% missing values. The identifiers and links to other sources are used only as gold standard.
\end{table}

\paragraph{Sources considered but not used.}
\emph{DBpedia} \citep{lehmann2015dbpedia} is extracted from the English Wikipedia, on which Wikidata
partly draws as well: 63\% of its linked pairs with Wikidata have coordinates identical to about ten metres,
so it is not distinct enough. \emph{OpenStreetMap} \citep{osm2026} (47,706 universities) is
independent but gives little beyond name and position (founding year for 2\%). The \emph{THE ranking}
\citep{the2025rankings} has only 3,118 institutions and shares only name, country and student numbers,
and the \emph{Hipo domain list} \citep{hipo2026domains} adds nothing that our sources lack.

\paragraph{Integrated schema.}
We applied the rules of thumb for schema integration to the attribute lists of Table~\ref{tab:datasets}.
The integrated schema has one class, \emph{University}, with the 21 attributes of
Table~\ref{tab:schema}; 11 are in at least two datasets, five besides the name are in all three
(alternative names, country, city, website, founding year) and seven are lists. Two mappings are not
one-to-one: ROR names a parent organisation by identifier, which we replace by its name, and the
\emph{types} of ROR are registry categories such as \emph{education} or \emph{funder} and are therefore
not mapped to the institution types of Wikidata. The internet domain is derived from the website for
Wikidata and DEQAR. No attribute is a key; the source identifiers are not part of the schema, and a
university is identified by name and country, which is unique for 99.1\% (Wikidata), 99.8\% (ROR)
and 94.5\% (DEQAR) of the records.

\begin{table}[t]
\caption{Attribute intersection with the integrated schema (W = Wikidata, R = ROR, D = DEQAR)}\label{tab:schema}
\scriptsize
\setlength{\tabcolsep}{3pt}
\begin{tabular}{@{}llc@{\hspace{6pt}}llc@{}}
\toprule
Attribute & Type & Found in & Attribute & Type & Found in\\
\midrule
name & string & W, R, D & parent organisation & string/list & R, D\\
alternative names & string/list & W, R, D & region, status & string & R\\
country & string & W, R, D & internet domains & string/list & R, W, D\\
city & string & W, R, D & institution types & string/list & W\\
website & string & W, R, D & number of students & integer & W\\
founding year & integer & W, R, D & memberships & string/list & W\\
acronyms & string/list & W, R & closure year & integer & D\\
country code & string & W, R & history & string/list & D\\
latitude, longitude & decimal & W, R & QA reports, last date & integer, date & D\\
\bottomrule
\end{tabular}
\end{table}

\paragraph{Schema mapping and translation.}
We first defined the target schema as an XSD (\url{data/target.xsd}): the elements \emph{id,
name, alternativeNames, acronyms, location, web, lifecycle, parentInstitutions, institutionTypes,
numberOfStudents, memberships, qualityAssurance} and \emph{status}. A first MapForce mapping for
Wikidata and DEQAR had flaws: ROR was not translated, the lists were split with stray spaces, the
institution types were not split, names and countries were lower-cased, and the DEQAR parent
institutions, which are needed for the hard non-matches, were missing. So that all three sources follow
exactly the same rules, the final files \url{data/Wikidata_Mapping.xml}, \url{data/ROR_Mapping.xml} and
\url{data/deqar_Mapping.xml} are written by one script (\url{translate_to_xml.py}) with the element
structure of the first mapping. They contain 22,697, 30,131 and 8,979 institutions, keep the original
spelling of names and countries, and validate against the XSD. All following steps read these XML files.

\section{Phase II: Identity Resolution}\label{sec:phase2}

\subsection{Gold standard}\label{sec:gold}

The sources link each other through identifier columns (Wikidata with ROR,
DEQAR with ROR); two records are linked across Wikidata and DEQAR if they point to the same ROR
record. We use the links that connect two records one-to-one as \emph{matches}: 13,020
Wikidata--ROR, 3,391 DEQAR--ROR and 2,144 Wikidata--DEQAR pairs. The links are only used for the gold
standard; the identifiers never enter the matching. They are incomplete, so we do not derive
non-matches from missing links. Table~\ref{tab:gold} shows the \emph{non-matches}:
\emph{hard} are candidate pairs of the most similar names in which both records are linked to a
different record of the other dataset (so neither can be the other's partner), \emph{random} a random draw
of such pairs, and three kinds of corner cases that the proposal singled out: a \emph{part of} an
institution (a ROR record with a parent) against the institution or against another part, a closed
DEQAR institution against the successor named in its history (\emph{predecessor}), and
\emph{same name} in different countries. The pairs are split by their left record into a development
part (tuning) and a test part (reported), so that no record is in both.

\begin{table}[t]
\caption{Gold standard: matches and non-matches by kind}\label{tab:gold}
\scriptsize
\setlength{\tabcolsep}{4pt}
\begin{tabular}{lrrrrrr}
\toprule
Pair & Matches & Hard & Random & Part of & Predecessor & Same name\\
\midrule
Wikidata--ROR & 13{,}020 & 13{,}020 & 6{,}510 & 2{,}154 & 0 & 229\\
DEQAR--ROR & 3{,}391 & 3{,}391 & 1{,}695 & 969 & 270 & 90\\
Wikidata--DEQAR & 2{,}144 & 2{,}144 & 1{,}072 & 237 & 190 & 80\\
\bottomrule
\end{tabular}
\end{table}


\subsection{Blocking}\label{sec:blocking}

Comparing all records would take 0.7 billion pairs for Wikidata--ROR. We compared the blockers of PyDI
\citep{pydi2026}, which implement the standard techniques \citep{christen2012matching}, by \emph{pair completeness} (share of the linked pairs that survive blocking) and
number of candidate pairs (Table~\ref{tab:blocking}). Equality of the normalised name finds only half to
three quarters of the matches, because 25--52\% of the linked pairs differ in name even after
normalisation (local-language names, abbreviations); a country-and-first-word key and sorted
neighbourhood reach 68--84\%. Token blocking on the words of the names produces hundreds of millions of
pairs because words such as \emph{university} are everywhere. Our best strategy scopes the tokens of the
name, alternative names and acronyms to the country (\texttt{DE\_heidelberg}), drops tokens that occur
in more than 150 records of a country, adds rare tokens without country for records whose country
differs or is missing, and adds all pairs that share an internet domain. It keeps 99.5--99.7\% of the
matches with 0.11--0.60 million candidate pairs (reduction ratio above 99.9\%).

\begin{table}[t]
\caption{Blocking: pair completeness (PC) and number of candidate pairs}\label{tab:blocking}
\scriptsize
\setlength{\tabcolsep}{4pt}
\begin{tabular}{lrrrrrr}
\toprule
Strategy & \multicolumn{2}{c}{Wikidata--ROR} & \multicolumn{2}{c}{DEQAR--ROR} & \multicolumn{2}{c}{Wikidata--DEQAR}\\ & PC (\%) & Cand. & PC (\%) & Cand. & PC (\%) & Cand.\\
\midrule
Standard: normalised name & 75.7 & 11k & 50.3 & 2k & 59.9 & 2k\\
Standard: country + first word & 83.8 & 872k & 68.2 & 148k & 75.7 & 101k\\
Sorted neighbourhood (window 10) & 84.1 & 241k & 68.0 & 101k & 75.9 & 95k\\
Token: words of the name & -- & 237M & -- & 91M & -- & 84M\\
Token: country-scoped name tokens & 99.2 & 597k & 99.1 & 193k & 98.6 & 110k\\
Token (scoped) + same web domain & 99.7 & 597k & 99.6 & 194k & 99.5 & 110k\\
\bottomrule
\end{tabular}
\end{table}


\subsection{Matching}\label{sec:matching}

\paragraph{Comparators.}
For every pair we compute 11 similarities with PyDI's comparators: the name with four string measures
(Jaro-Winkler, Levenshtein, Jaccard on words, Sørensen-Dice), the best match over all names of the two
records (name, alternative names, acronyms) with Jaro-Winkler and with Jaccard, the country (1 if equal, 0.5 if unknown, 0 if different), the
shared internet domain, the city (Jaro-Winkler), the founding year (absolute difference scaled to 10
years) and the distance between the coordinates (1 up to 2\,km, 0 from 50\,km). A missing value gives 0.
Four indicator features tell whether both records have a domain, city, year or position.

\paragraph{Experiments.}
We compared (i) single similarity measures on the name, (ii) five rule-based matchers that are the
weighted average of an increasing set of attributes (PyDI's \texttt{RuleBasedMatcher}), with integer
weights from 0 to 6 found by random search on the development part, and (iii) three learned matchers
(logistic regression, random forest, gradient boosting of scikit-learn \citep{pedregosa2011sklearn}, in
PyDI's \texttt{MLBasedMatcher}) on all 15 features. Thresholds were tuned on the development part (learned
matchers: five-fold cross validation). Table~\ref{tab:matching} shows precision (P), recall (R) and F1
on the test part. As in PyDI's evaluator, a predicted pair is a false positive only if the gold standard
labels it as a non-match; matches lost by blocking count as missed.

\begin{table}[t]
\caption{Matching on the test part of the gold standard (precision, recall, F1 in \%); thresholds and weights tuned on the development part}\label{tab:matching}
\scriptsize
\setlength{\tabcolsep}{4pt}
\begin{tabular}{lrrrrrrrrr}
\toprule
Matcher & \multicolumn{3}{c}{Wikidata--ROR} & \multicolumn{3}{c}{DEQAR--ROR} & \multicolumn{3}{c}{Wikidata--DEQAR}\\ & P & R & F1 & P & R & F1 & P & R & F1\\
\midrule
Name, Jaro-Winkler & 96.4 & 77.0 & 85.6 & 90.2 & 54.7 & 68.1 & 93.3 & 65.7 & 77.1\\
Name, Levenshtein & 97.4 & 76.7 & 85.8 & 91.6 & 54.4 & 68.3 & 96.6 & 62.4 & 75.8\\
Name, Jaccard & 94.5 & 77.8 & 85.3 & 95.2 & 55.0 & 69.7 & 96.0 & 66.7 & 78.7\\
Name, Sørensen-Dice & 94.5 & 77.8 & 85.3 & 95.2 & 55.0 & 69.7 & 96.0 & 66.7 & 78.7\\
All names, best match, Jaro-Winkler & 93.2 & 95.2 & 94.2 & 95.9 & 80.8 & 87.7 & 95.3 & 82.5 & 88.4\\
All names, best match, Jaccard & 93.5 & 95.0 & 94.3 & 95.8 & 82.3 & 88.5 & 95.5 & 83.2 & 88.9\\
\midrule
names & 93.2 & 95.2 & 94.2 & 91.2 & 82.5 & 86.6 & 95.3 & 82.5 & 88.4\\
names+country & 98.0 & 94.4 & 96.2 & 95.8 & 80.6 & 87.6 & 95.3 & 82.5 & 88.4\\
names+country+domain & 95.9 & 98.1 & 97.0 & 97.7 & 93.3 & 95.4 & 96.8 & 94.7 & 95.7\\
names+country+domain+city & 95.8 & 98.1 & 97.0 & 96.4 & 92.9 & 94.6 & 96.8 & 94.7 & 95.7\\
all attributes & 96.5 & 93.0 & 94.7 & 98.3 & 90.1 & 94.1 & 96.2 & 94.7 & 95.4\\
\midrule
logistic regression & 97.9 & 97.3 & 97.6 & 97.7 & 93.6 & 95.6 & 97.4 & 94.9 & 96.1\\
random forest & 99.1 & 98.2 & 98.6 & 98.9 & 93.6 & 96.2 & 96.6 & 96.1 & 96.4\\
gradient boosting & 99.5 & 97.6 & 98.6 & 98.3 & 94.2 & 96.2 & 98.1 & 94.1 & 96.1\\
\bottomrule
\end{tabular}
\end{table}


\paragraph{Findings.}
(1)~The choice of the string measure hardly matters (F1 85.3--85.8\% for Wikidata--ROR), but comparing
\emph{all} names of two records instead of the main names raises F1 by 8--19 points, because the
main names often differ by language. (2)~The web domain is the most useful further attribute (+0.8 to
+7.8 points on top of names and country); the country helps most for Wikidata--ROR (+2.0 points), where
names such as ``Victoria University'' exist in several countries. (3)~Adding all attributes to a
weighted average does \emph{not} help (``all attributes'' is 0.3--2.3 points below the rule with names,
country and domain): the founding year and the coordinates are missing for many records, and a missing
value counts as 0 and lowers the score of true matches. (4)~Learned matchers, which see the indicator
features, are best by 0.7--1.6 points over the best rule: random forest for Wikidata--ROR and gradient
boosting for DEQAR--ROR and Wikidata--DEQAR were chosen by their development F1. The domain is the most important
feature of the DEQAR pairs (importance 0.36 and 0.80), then the best name match. Recall at the final
threshold (0.8, see below) is 97\%, 93\% and 95\% for test matches in which both records have a website but
86\%, 64\% and 47\% when one has none (414, 25 and 19 matches). (5)~The final threshold misses 3.5\%, 7.4\% and 5.6\% of the
test matches. They are mostly records with names in different languages and different
websites (e.g.\ \emph{Academy of fine arts of Roma} -- \emph{Accademia di Belle Arti di Roma}) or a
campus against its university (\emph{Oklahoma State University} -- \emph{Oklahoma State
University-Main Campus}); 5--7\% of the missed matches were already lost by blocking.

\subsection{One-to-one matching and the final threshold}

Every institution should be matched at most once per dataset, so we compared PyDI's post-clustering
algorithms (Table~\ref{tab:post}). Greedy and stable one-to-one matching give the same
result and raise precision to 99.1--100\% at a recall loss of 0.6--1.0 points; maximum bipartite
matching is worse. We use greedy matching.

The F1-optimal thresholds were found on a gold standard whose non-matches are mostly easy. We therefore
checked, for each dataset pair, 60 randomly drawn correspondences that are \emph{not} in the gold standard
by hand (Table~\ref{tab:threshold}, columns ``Sample''). At the F1-optimal threshold only 75\% of
the new Wikidata--ROR pairs were correct, the errors being similar institutions in one city
(\emph{British University of Nicosia} -- \emph{Cyprus Aydin University}), a part against the whole
(\emph{Kmct Polytechnic College} -- \emph{KMCT Group of Educational Institutions}) and predecessors
(\emph{University of Auvergne} -- \emph{Universit\'e Blaise Pascal}). A threshold of 0.8 raises the
precision of the new pairs to 97\% (Wikidata--ROR), 94\% (DEQAR--ROR) and 98\% (Wikidata--DEQAR) for a
loss of at most 1.6 points of recall on the gold standard. We use 0.8 for all pairs. The resulting
20,079 correspondences contain 2,318 pairs that are not in the gold standard, i.e.\ links that the
sources do not give; the sample suggests that nearly all of them are correct.

\begin{table}[t]
\caption{Post-clustering of the correspondences (threshold of the matcher, test part of the gold standard)}\label{tab:post}
\scriptsize
\setlength{\tabcolsep}{4pt}
\begin{tabular}{llrrrr}
\toprule
Pair & Post-clustering & Correspondences & P & R & F1\\
\midrule
Wikidata--ROR & no post-clustering & 14{,}625 & 99.1 & 98.2 & 98.6\\
 & Greedy one-to-one & 13{,}293 & 100.0 & 97.6 & 98.8\\
 & Maximum bipartite matching & 13{,}301 & 99.7 & 95.1 & 97.4\\
 & Stable matching & 13{,}290 & 100.0 & 97.6 & 98.8\\
DEQAR--ROR & no post-clustering & 4{,}818 & 98.3 & 94.2 & 96.2\\
 & Greedy one-to-one & 4{,}278 & 99.6 & 93.6 & 96.5\\
 & Maximum bipartite matching & 4{,}281 & 99.1 & 90.7 & 94.7\\
 & Stable matching & 4{,}277 & 99.6 & 93.6 & 96.5\\
Wikidata--DEQAR & no post-clustering & 3{,}666 & 98.1 & 94.1 & 96.1\\
 & Greedy one-to-one & 3{,}010 & 99.1 & 93.1 & 96.0\\
 & Maximum bipartite matching & 3{,}014 & 99.0 & 91.4 & 95.0\\
 & Stable matching & 3{,}008 & 99.1 & 93.1 & 96.0\\
\bottomrule
\end{tabular}
\end{table}

\begin{table}[t]
\caption{Threshold analysis after greedy one-to-one matching: gold standard (test part) and hand-checked sample of the correspondences that are not in the gold standard}\label{tab:threshold}
\scriptsize
\setlength{\tabcolsep}{4pt}
\begin{tabular}{lrrrrrrr}
\toprule
Pair & Threshold & Corr. & P & R & F1 & Sample & Sample precision\\
\midrule
Wikidata--ROR & 0.5 & 13{,}293 & 100.0 & 97.6 & 98.8 & 60 & 75.0\\
 & 0.7 & 13{,}048 & 100.0 & 96.8 & 98.4 & 42 & 92.9\\
 & 0.8 & 12{,}912 & 100.0 & 96.0 & 98.0 & 39 & 97.4\\
 & 0.9 & 12{,}609 & 100.0 & 93.7 & 96.8 & 32 & 100.0\\
DEQAR--ROR & 0.5 & 4{,}278 & 99.6 & 93.6 & 96.5 & 60 & 88.3\\
 & 0.7 & 4{,}231 & 99.7 & 93.0 & 96.2 & 58 & 91.4\\
 & 0.8 & 4{,}157 & 99.7 & 92.0 & 95.7 & 54 & 94.4\\
 & 0.9 & 4{,}039 & 99.7 & 89.8 & 94.5 & 49 & 93.9\\
Wikidata--DEQAR & 0.5 & 3{,}010 & 99.1 & 93.1 & 96.0 & 60 & 98.3\\
 & 0.7 & 3{,}010 & 99.1 & 93.1 & 96.0 & 60 & 98.3\\
 & 0.8 & 3{,}010 & 99.1 & 93.1 & 96.0 & 60 & 98.3\\
 & 0.9 & 2{,}929 & 99.3 & 92.5 & 95.8 & 57 & 98.2\\
\bottomrule
\end{tabular}
\end{table}


\section{Phase III: Data Fusion}\label{sec:phase3}

\paragraph{Groups.}
We merge the three correspondence sets into groups (connected components) with PyDI. Because the
sets are matched pairwise, they can contradict each other (two Wikidata records linked to one
ROR record through DEQAR); we therefore accept correspondences from the highest score down and reject
one that would put two records of the same dataset into a group. This rejected 49 of 20,079
correspondences. The 61,807 records form 44,340 groups, i.e.\ \textbf{44,340 entities}: 15,771 only in ROR,
9,465 only in Wikidata, 4,474 only in DEQAR, 10,125 in ROR and Wikidata, 1,398 in DEQAR and ROR, 270 in
DEQAR and Wikidata and 2,837 in all three.

\paragraph{Strategies.}
The conflict resolution functions follow the usual taxonomy \citep{bleiholder2009fusion}. For the list attributes (alternative names, acronyms, domains, parent institutions, types, memberships,
history) we take the union. For the single-valued attributes we compared seven strategies with PyDI's
fusion engine: PyDI's default (first non-null value), \emph{favour a source} with each of three orders
(ROR, Wikidata or DEQAR first; PyDI's \texttt{prefer\_higher\_trust}), and majority \emph{voting} on
normalised values (own resolver) with ties decided in favour of ROR, in favour of Wikidata, or of the
earliest year. Coordinates are fused with the median of the values. Attributes of only one source need no
resolution. In 99.4\% of the groups with two country codes the sources agree, but only 70\% agree on
the name, 74\% on the city, 80\% on the founding year and 86\% on the web domain; the rest are the conflicts
that fusion has to resolve. With two sources most votes are ties, so a vote is in effect a
trust-based choice among the disagreeing pairs.

\paragraph{Evaluation.}
We evaluated the strategies on a fusion gold standard of 39 well-known universities that are in all
three datasets (\url{data/gold_standard/fusion_gold.csv}); the expected name (English or official
local name), country, city, founding year and web domain were compiled by hand from general knowledge
of these institutions, not from one of the sources. Table~\ref{tab:fusion} shows the accuracy. Voting is
as good as or better than every rule that favours a source, except for the founding year. Favouring
DEQAR, which is also what PyDI's default does here, is least accurate (founding year 89.7\%, name 92.3\%:
three official local names are not among the accepted names); favouring Wikidata sometimes takes a district as
the city (Wikidata gives ``Taka-T\"o\"ol\"o'' for the Sibelius Academy in Helsinki) and is wrong for 13\% of the entities. For the
founding year, Wikidata first (97.4\%) is one entity better than voting (94.9\%), which is not a
reliable difference on 39 entities. Our final strategy takes, per attribute,
the best accuracy and decides ties in favour of voting: voting for name, country, city and domain, and
Wikidata first for the founding year. The set is small and its universities are well documented, so
it can show that a strategy is bad but not rank the good ones.

\begin{table}[t]
\caption{Accuracy (\%) of the fusion strategies on the 39 entities of the fusion gold standard}\label{tab:fusion}
\scriptsize
\setlength{\tabcolsep}{4pt}
\begin{tabular}{lrrrrr}
\toprule
Strategy & Name & Country & City & Founding year & Domain\\
\midrule
first non-null (PyDI default) & 92.3 & 100.0 & 100.0 & 89.7 & 97.4\\
ROR first & 100.0 & 100.0 & 100.0 & 94.9 & 97.4\\
Wikidata first & 100.0 & 100.0 & 87.2 & 97.4 & 100.0\\
DEQAR first & 92.3 & 100.0 & 100.0 & 89.7 & 97.4\\
vote, tie: ROR & 100.0 & 100.0 & 100.0 & 94.9 & 100.0\\
vote, tie: Wikidata & 100.0 & 100.0 & 100.0 & 94.9 & 100.0\\
vote, tie: earliest year & 100.0 & 100.0 & 100.0 & 94.9 & 100.0\\
\bottomrule
\end{tabular}
\end{table}


\paragraph{Quality of the fused dataset.}
Fusion fills the gaps of the sources (Table~\ref{tab:completeness}). For the entities that are in at
least two sources the fused record has a founding year in 99\%, a city in 100\%, coordinates in
99\% and alternative names in 91\%, where the best single source has 94\%, 98\%, 98\% and 76\%.
Over all 44,340 entities, attributes that only one source provides remain sparse
(students 8\%, memberships 11\%, quality assurance reports 11\%) because
only few of their entities are in that source.

\begin{table}[t]
\caption{Share of entities with a value (\%): source datasets and fused dataset over all records, and best single source versus fused record for the entities that are in at least two sources}\label{tab:completeness}
\scriptsize
\setlength{\tabcolsep}{4pt}
\begin{tabular}{lrrrrrr}
\toprule
Attribute & Wikidata & ROR & DEQAR & Fused & Best source & Fused \\ & \multicolumn{4}{c}{all entities} & \multicolumn{2}{c}{entities in $\geq 2$ sources}\\
\midrule
Alternative names & 53 & 71 & 69 & 66 & 76 & 91\\
Acronyms & 12 & 54 & 0 & 40 & 60 & 65\\
Country code & 95 & 100 & 99 & 98 & 98 & 100\\
City & 83 & 100 & 74 & 89 & 98 & 100\\
Region & 0 & 100 & 0 & 68 & 98 & 98\\
Web domain & 75 & 97 & 88 & 85 & 97 & 99\\
Founding year & 81 & 89 & 59 & 80 & 94 & 99\\
Coordinates & 69 & 100 & 0 & 79 & 98 & 99\\
Institution types & 100 & 0 & 0 & 51 & 90 & 90\\
Number of students & 16 & 0 & 0 & 8 & 23 & 23\\
Memberships & 22 & 0 & 0 & 11 & 33 & 33\\
Parent institutions & 0 & 9 & 28 & 12 & 7 & 12\\
QA reports & 0 & 0 & 57 & 11 & 22 & 22\\
\bottomrule
\end{tabular}
\end{table}


\section{Conclusion}

All requirements of the course are met: three datasets from distinct sources, 44,340 entities
(61,807 records), 14,630 entities in at least two and 2,837 in all three datasets, 21 attributes of
which five besides the name are in all three, and seven list attributes. The main lessons are: names alone
are not enough across languages (the best match over all names and the web domain matter most);
learned matchers beat hand-weighted rules mainly because they handle missing values; a gold standard
whose non-matches are easy gives thresholds that are too low, which a hand-checked sample of the
correspondences revealed; and with only two sources voting is mostly a trust decision.

\begin{table}[t]
\caption{Project requirements (slide 12) and how the project meets them}\label{tab:requirements}
\scriptsize
\setlength{\tabcolsep}{3pt}
\begin{tabularx}{\textwidth}{@{}lP{6.2cm}L@{}}
\toprule
& Requirement & Our project\
\midrule
1. & 3 different data sets & 3: Wikidata, ROR, DEQAR\
2. & at least 2,500 entities in total (good: $>$\,10,000 but $<$\,100,000) & 44,340 entities (61,807 records)\
3. & at least 1,000 entities in at least two datasets & 14,630 (2,837 in all three)\
4. & at least 8 attributes; entities identifiable by at least two attributes & 21; name + country\
5. & at least 5 attributes in at least two datasets; some in three & 11; 5 other than the name in all three\
6. & ideally at least one list attribute & 7, e.g.\ alternative names, history\
\bottomrule
\end{tabularx}
\end{table}

\paragraph{Limitations.}
The matches of the gold standard are the pairs that the sources link themselves, so they are
over-represented among records with a website and a Wikidata or ROR identifier, and the
recall on unlinked records is unknown. The hand-checked samples (60 pairs per dataset pair) give a
precision with a margin of roughly $\pm 6$ points, and the fusion gold standard has 39 entities.
Parts and predecessors are kept as separate entities but the data cannot tell a branch from a
university in general, so non-universities of ROR and DEQAR remain in the fused dataset.

\bibliographystyle{plainnat}
\bibliography{references}

\end{document}
